\documentclass[../../main.tex]{subfiles}

\begin{document}

\chapter{Interupsi BIOS/DOS dan Operasi I/O}
\label{chap:interupsi}

% =========================================================
% OUTCOME PEMBELAJARAN KHUSUS BAB
% =========================================================

\begin{subcpmkbox}
\SubCPMK{2.4} --- Menjelaskan mekanisme interupsi dan layanan BIOS/DOS
(\code{INT 10h}, \code{INT 21h}) sebagai jembatan antara instruksi mesin
dan perangkat keras serta operasi masukan/keluaran.

\vspace{0.4em}
\textbf{CPMK yang didukung:} \CPMK{02} \quad\textbullet\quad
\textbf{CPL yang didukung:} \CPL{02}
\end{subcpmkbox}

\begin{tujuanbox}
Setelah mempelajari bab ini, mahasiswa mampu:
\begin{enumerate}
    \item menjelaskan alasan diperlukannya interupsi dan membedakan
          interupsi perangkat keras, perangkat lunak, serta
          pengecualian prosesor;
    \item menghitung alamat vektor interupsi ke-$n$ dan menyebutkan
          vektor-vektor penting beserta kegunaannya;
    \item menelusuri urutan penyimpanan \reg{FLAGS}, \reg{CS}, dan
          \reg{IP} pada stack serta pengambilannya kembali oleh
          \asm{IRET};
    \item menggunakan layanan BIOS \code{INT 10h}, \code{INT 16h},
          \code{INT 17h} serta layanan DOS \code{INT 21h};
    \item memasang \eng{interrupt handler} sederhana dengan mengganti
          isi sebuah vektor, lalu memulihkan vektor tersebut.
\end{enumerate}
\end{tujuanbox}

\begin{materibox}
\begin{enumerate}
    \item Konsep interupsi: perangkat keras, perangkat lunak, dan
          pengecualian.
    \item Tabel vektor interupsi pada alamat \hex{0000}--\hex{03FF}
          dan cara menghitung alamat vektor ke-$n$.
    \item Mekanisme \asm{INT} dan \asm{IRET} serta kaitannya dengan
          \flag{TF} dan eksekusi satu langkah.
    \item Layanan BIOS: \code{INT 10h}, \code{INT 16h}, dan
          \code{INT 17h}.
    \item Layanan DOS \code{INT 21h} beserta penyangga fungsi
          \hex{0A} dan perbedaan \lstasm|INT 20h| dengan fungsi
          \hex{4C}.
    \item Penanganan interupsi buatan sendiri beserta risikonya, serta
          studi kasus program sapaan dan teks berwarna.
\end{enumerate}
\end{materibox}

\subfile{section-01}
\subfile{section-02}
\subfile{section-03}

\section{Contoh}
\label{sec:contoh-bab14}

Dua contoh berikut memperlihatkan prosedur berparameter dan macro.
Contoh pertama mengirim parameter melalui stack. Contoh kedua
memakai macro dengan parameter karakter.

\subsection{Contoh 1: Prosedur Berparameter melalui Stack}

Program pada Kode~\ref{lst:prosedur-luas} menghitung luas persegi
panjang melalui prosedur \code{HitungLuas}.

\lstinputlisting[
    language=MASM,
    caption={Program \file{prosedur-luas.asm} memakai prosedur
    berparameter melalui stack.},
    label={lst:prosedur-luas}
]{\CodeDir/prosedur-luas.asm}

Bagian penting program:

\begin{itemize}
    \item Pemanggil mendorong \code{tinggi} lebih dahulu, lalu
          \code{lebar}. Parameter yang didorong terakhir berada paling
          dekat dengan puncak stack.
    \item \asm{CALL HitungLuas} menyimpan alamat kembali ke stack.
    \item Prosedur menyimpan \reg{BP} lama dengan \asm{PUSH BP}, lalu
          menyalin \reg{SP} ke \reg{BP}.
    \item \lstasm|MOV AX, [BP+4]| membaca parameter pertama, yaitu
          \code{lebar}.
    \item \lstasm|MUL WORD PTR [BP+6]| mengalikan dengan parameter
          kedua, yaitu \code{tinggi}.
    \item \asm{RET} \code{4} kembali ke pemanggil sekaligus membuang
          empat byte parameter.
\end{itemize}

Hasilnya adalah $5 \times 10 = 50$ atau \hex{0032}, disimpan pada
variabel \code{luas} dan register \reg{AX}.

\subsection{Contoh 2: Macro dengan Parameter}

Program pada Kode~\ref{lst:macro-demo} mendefinisikan macro
\code{Cetak\_Kar} yang mencetak satu karakter sebanyak tiga kali.

\begin{lstlisting}[
    language=MASM,
    caption={Macro \code{Cetak\_Kar} dan pemakaiannya.},
    label={lst:macro-demo}
]
.MODEL SMALL
.STACK 100h

Cetak_Kar MACRO Kar
    MOV CX, 3
    MOV AH, 02h
    MOV DL, Kar
Ulang:
    INT 21h
    LOOP Ulang
ENDM

.CODE
START:
    Cetak_Kar 'S'       ; cetak 'S' tiga kali

    MOV AH, 4Ch
    MOV AL, 00h
    INT 21h
END START
\end{lstlisting}

Saat assembler menemukan \lstasm|Cetak_Kar 'S'|, isi macro disalin ke
tempat itu dengan setiap kemunculan \code{Kar} diganti \code{'S'}.
Keluaran program adalah \code{SSS}.

\begin{lstlisting}[
    style=sessiongaya,
    caption={Keluaran program macro.},
    label={lst:keluaran-macro}
]
SSS
\end{lstlisting}

\begin{infobox}
\vspace{1.1em}
Macro pada contoh ini memakai label \code{Ulang}. Label tersebut
aman karena macro hanya dipakai sekali. Bila macro dipakai lebih dari
sekali, label yang sama akan disalin berulang dan assembler dapat
melaporkan kesalahan label ganda.
\end{infobox}

\section{Praktikum}
\label{sec:praktikum-bab14}

\begin{praktikumbox}
\textbf{Judul:} Operasi masukan/keluaran dengan \code{INT 21h} dan
penelusuran stack pada pemanggilan \asm{INT}.\\[0.3em]
\textbf{Sub-CPMK:} 2.4 \quad\textbullet\quad
\textbf{Alokasi waktu:} 100 menit \quad\textbullet\quad
\textbf{Bentuk:} individu \quad\textbullet\quad
\textbf{Tugas 7 dikumpulkan pada praktikum ini.}
\end{praktikumbox}

\subsection{Tujuan Praktikum}

Setelah menyelesaikan praktikum ini, mahasiswa mampu menulis program
masukan/keluaran dengan \code{INT 21h} dan menunjukkan isi stack
sebelum serta sesudah instruksi \asm{INT} dijalankan.

\subsection{Alat dan Bahan}

\begin{itemize}
    \item Komputer dengan Emu8086, GUI Turbo Assembler, atau DOSBox
          beserta DEBUG.
    \item Berkas \file{input-nama.asm} dan \file{warna-layar.asm} dari
          direktori \file{code/asm/}.
    \item Jurnal praktikum untuk mencatat nilai register dan isi
          stack.
\end{itemize}

\subsection{Langkah Kerja}

\subsubsection{A. Masukan/keluaran dengan \code{INT 21h}}

\begin{enumerate}
    \item Rakit dan jalankan \file{input-nama.asm} seperti pada
          Kode~\ref{lst:input-nama}, lalu catat keluarannya.
    \item Ubah kapasitas penyangga dari \dec{20} menjadi \dec{5}, lalu
          ketik nama yang lebih panjang dari empat karakter dan
          jelaskan hasilnya dengan mengacu pada
          Kode~\ref{lst:buffer-0a}.
    \item Hapus baris \lstasm|MOV BYTE PTR [SI], '$'|, rakit ulang, dan
          jalankan. Amati keluarannya dan jelaskan penyebabnya.
    \item Tambahkan satu baris keluaran yang mencetak NIM mahasiswa
          memakai fungsi \hex{09}.
\end{enumerate}

\subsubsection{B. Menelusuri Stack pada Pemanggilan \asm{INT}}

\begin{enumerate}
    \item Buka DOSBox, lakukan \file{mount} ke direktori kerja, lalu
          rakit program dengan \file{tasm} dan \file{tlink}.
    \item Jalankan \file{debug input-nama.exe} dan tampilkan instruksi
          program dengan perintah \file{U}. Tandai alamat instruksi
          \lstasm|INT 21h| yang memanggil fungsi \hex{09}.
    \item Tampilkan register dengan perintah \file{R}, lalu catat
          nilai \reg{SP}, \reg{CS}, \reg{IP}, dan \reg{FLAGS}.
    \item Jalankan \file{T} sampai tepat sebelum \lstasm|INT 21h|,
          catat nilai \reg{SP}, lalu tampilkan puncak stack dengan
          \file{D SS:SP L 6} dan catat enam byte yang tampak.
    \item Jalankan \file{T} satu kali untuk mengeksekusi
          \lstasm|INT 21h|, lalu ulangi perintah \file{D SS:SP L 6}.
    \item Bandingkan hasil langkah 4 dan 5, jelaskan mengapa \reg{SP}
          berubah, dan tentukan nilai mana yang merupakan \reg{IP},
          \reg{CS}, serta \reg{FLAGS}.
    \item Ulangi langkah 4 dan 5 dengan \file{P} sebagai pembanding,
          lalu bandingkan hasilnya dengan jendela \file{Stack} pada
          Emu8086.
\end{enumerate}

\begin{infobox}
Perintah \file{T} (\eng{trace}) menjalankan satu instruksi dan
\emph{masuk} ke rutin penangan, sedangkan \file{P} (\eng{proceed})
menjalankannya sampai selesai tanpa masuk ke penangan. Keduanya
memanfaatkan \flag{TF} seperti dijelaskan pada
Subbab~\ref{sec:konsep-interupsi}.
\end{infobox}

\subsubsection{C. Tugas 7: Sapaan Berwarna}

Tugas 7 menggabungkan kedua contoh program pada bab ini. Buatlah
program baru bernama \file{tugas7.asm} dengan ketentuan berikut.

\begin{enumerate}
    \item Program meminta nama pengguna dengan \code{INT 21h} fungsi
          \hex{0A} dan kapasitas penyangga \dec{30} karakter, lalu
          meminta NIM pada penyangga terpisah.
    \item Program menetapkan mode video \hex{03} dengan
          \code{INT 10h} fungsi \hex{00}, lalu mencetak baris
          \code{Halo, <nama>} dengan fungsi \hex{13} dan atribut
          kuning.
    \item Program mencetak baris \code{NIM: <nim>} dengan atribut
          hijau terang, lalu menutup tampilan dengan garis merah
          seperti pada \file{warna-layar.asm}.
    \item Program diakhiri \code{INT 21h} fungsi \hex{4C} dan memuat
          komentar berbahasa Indonesia pada setiap baris penting.
\end{enumerate}

\subsection{Hasil yang Dikumpulkan}

\begin{enumerate}
    \item Berkas \file{tugas7.asm} beserta hasil perakitan dan
          tangkapan layar keluarannya.
    \item Tangkapan layar keluaran program pada langkah A.3 beserta
          penjelasan penyebabnya.
    \item Tabel pengamatan nilai \reg{SP} dan isi puncak stack sebelum
          serta sesudah \lstasm|INT 21h|, beserta jawaban langkah B.6
          dan B.7.
\end{enumerate}

\begin{asesmenbox}
\textbf{Asesmen praktikum Bab~\ref{chap:interupsi} (bobot 10\% dari
komponen praktikum).}

\begin{tabular}{@{}p{5.9cm}cp{5.0cm}@{}}
    \toprule
    \textbf{Indikator} & \textbf{Bobot} & \textbf{Kriteria} \\
    \midrule
    Program masukan/keluaran dengan \code{INT 21h} berjalan benar &
    20\% & Ketiga fungsi (\hex{0A}, \hex{09}, \hex{02}) dipakai
    dengan argumen yang tepat \\
    \addlinespace
    Tugas 7 selesai sesuai ketentuan & 35\% & Seluruh ketentuan
    terpenuhi dan teks berwarna tampil benar \\
    \addlinespace
    Ketepatan tabel pengamatan stack & 25\% & Nilai \reg{SP} dan isi
    stack tercatat benar pada setiap langkah \\
    \addlinespace
    Penjelasan hasil pengamatan & 10\% & Alasan perubahan \reg{SP}
    dan makna \file{T} serta \file{P} dijelaskan dengan tepat \\
    \addlinespace
    Kerapian kode dan jurnal praktikum & 10\% & Komentar lengkap dan
    jurnal dapat diulang orang lain \\
    \bottomrule
\end{tabular}
\end{asesmenbox}

\section{Latihan}
\label{sec:latihan-bab14}

\begin{enumerate}
    \item \textbf{Stack.} Jelaskan pengertian \eng{last in, first
          out}. Mengapa stack 8086 tumbuh ke alamat yang lebih
          rendah?

    \item \textbf{PUSH dan POP.} Diketahui \reg{SP} = \hex{FFFE}.
          Tentukan nilai \reg{SP} setelah \asm{PUSH AX} dan setelah
          \asm{POP BX}. Jelaskan urutan langkah masing-masing.

    \item \textbf{Prosedur.} Jelaskan fungsi \asm{PROC}, \asm{ENDP},
          \asm{CALL}, dan \asm{RET}. Mengapa \asm{CALL} berbeda dari
          \asm{JMP}?

    \item \textbf{Stack frame.} Diketahui prosedur menerima dua
          parameter melalui stack. Jelaskan isi \code{[BP+0]},
          \code{[BP+2]}, \code{[BP+4]}, dan \code{[BP+6]} setelah
          \asm{PUSH BP} dan \asm{MOV BP, SP}.

    \item \textbf{Penyampaian data.} Sebutkan tiga cara menyampaikan
          data ke prosedur. Berikan kelebihan dan kekurangan
          masing-masing.

    \item \textbf{RET n.} Jelaskan fungsi \asm{RET} \code{4}.
          Kapan pembersihan stack dilakukan oleh prosedur, dan kapan
          oleh pemanggil?

    \item \textbf{Macro.} Jelaskan cara mendefinisikan macro dengan
          parameter. Apa yang terjadi pada saat assembler menemukan
          pemakaian macro?

    \item \textbf{Perbandingan.} Sebutkan empat perbedaan prosedur
          dan macro. Kapan sebaiknya memakai prosedur, dan kapan
          memakai macro?
\end{enumerate}

\begin{refleksibox}
\vspace{1.1em}
Jawab pertanyaan berikut pada jurnal praktikum Anda.
\begin{enumerate}
    \item Bagian mana yang paling sulit Anda telusuri: perubahan
          \reg{SP}, pembentukan \eng{stack frame}, atau penyalinan
          macro? Mengapa?
    \item Setelah memeriksa stack dengan DEBUG, apakah Anda memahami
          hubungan antara \asm{PUSH}, \asm{CALL}, dan \asm{RET}?
          Jelaskan.
    \item Konsep apa yang perlu Anda ulangi sebelum mempelajari
          array, string, dan pencetakan bilangan pada bab berikutnya?
\end{enumerate}
\end{refleksibox}

\section{Rangkuman}
\label{sec:rangkuman-bab14}

Pada bab ini telah dibahas:

\begin{itemize}
    \item Stack adalah daerah memori dengan aturan \eng{last in,
          first out}. Stack 8086 ditunjuk \reg{SS}:\reg{SP} dan
          tumbuh ke alamat yang lebih rendah.

    \item \asm{PUSH} mengurangi \reg{SP} dua dan menyimpan satu word.
          \asm{POP} mengambil satu word dan menambah \reg{SP} dua.

    \item Prosedur didefinisikan dengan \asm{PROC} dan \asm{ENDP}.
          Pemanggil memakai \asm{CALL}, sedangkan prosedur kembali
          dengan \asm{RET}.

    \item \asm{CALL} menyimpan alamat kembali ke stack. \asm{RET}
          mengambil alamat itu dan memuatnya kembali ke \reg{IP}.

    \item \eng{Stack frame} dibentuk dengan \asm{PUSH BP} dan
          \asm{MOV BP, SP}. Parameter pertama berada pada
          \code{[BP+4]}, parameter kedua pada \code{[BP+6]}.

    \item Data dapat disampaikan ke prosedur melalui register,
          \eng{stack}, atau variabel. \asm{RET} \code{n} membuang
          $n$ byte parameter setelah kembali.

    \item Macro didefinisikan dengan \asm{MACRO} dan diakhiri
          \asm{ENDM}. Macro dapat memiliki parameter dan disalin ke
          tempat pemakaian pada saat perakitan.

    \item Prosedur menekankan ukuran program dan modularitas,
          sedangkan macro menekankan kemudahan parameter dan
          menghilangkan biaya \asm{CALL}/\asm{RET}.
\end{itemize}

\begin{checklistbox}
\vspace{1.1em}
Tandai setiap butir yang sudah Anda kuasai.

\begin{itemize}
    \item[\checkbox] Saya dapat menjelaskan sifat LIFO pada stack.
    \item[\checkbox] Saya dapat menggunakan \asm{PUSH} dan
          \asm{POP}.
    \item[\checkbox] Saya dapat menjelaskan perubahan \reg{SP} pada
          operasi stack.
    \item[\checkbox] Saya dapat menulis prosedur dengan
          \asm{PROC}/\asm{ENDP}.
    \item[\checkbox] Saya dapat menggunakan \asm{CALL} dan
          \asm{RET}.
    \item[\checkbox] Saya dapat mengakses parameter melalui
          \eng{stack frame} dan \reg{BP}.
    \item[\checkbox] Saya dapat mendefinisikan dan memakai macro
          berparameter.
    \item[\checkbox] Saya dapat membedakan prosedur dan macro.
\end{itemize}

\vspace{0.4em}
Bila masih ada butir yang belum bertanda, ulangi bagian yang
bersesuaian sebelum melanjutkan ke bab berikutnya.
\end{checklistbox}

\begin{infobox}
\vspace{1.1em}
\textbf{Rujukan bab.} Materi pada bab ini diringkas dari
\cite{gunawan2019}, Bagian 12--13, hlm.~115--126, dan
\cite{sto2001}, BAB XIV, hlm.~73--77, serta BAB XVI--XVII,
hlm.~85--92.
\end{infobox}

\section{Evaluasi}
\label{sec:evaluasi-bab14}

\subsection{Soal Pilihan Ganda}

\begin{enumerate}
    \item Setelah \asm{PUSH AX} dijalankan, nilai \reg{SP} \dots
          \begin{enumerate}
              \item bertambah 2
              \item berkurang 2
              \item tidak berubah
              \item bertambah 4
          \end{enumerate}

    \item Alamat kembali pada pemanggilan prosedur disimpan oleh
          instruksi \dots
          \begin{enumerate}
              \item \asm{JMP} \item \asm{CALL}
              \item \asm{RET} \item \asm{PUSH}
          \end{enumerate}

    \item Setelah \asm{PUSH BP} dan \asm{MOV BP, SP}, parameter
          pertama pada stack frame diakses melalui \dots
          \begin{enumerate}
              \item \code{[BP+0]} \item \code{[BP+2]}
              \item \code{[BP+4]} \item \code{[BP+6]}
          \end{enumerate}

    \item Macro diakhiri dengan direktif \dots
          \begin{enumerate}
              \item \asm{ENDP} \item \asm{ENDM}
              \item \asm{ENDS} \item \asm{END}
          \end{enumerate}

    \item Perbedaan utama prosedur dan macro adalah \dots
          \begin{enumerate}
              \item prosedur memakai \asm{CALL}, macro disalin saat
                    perakitan
              \item prosedur tidak dapat memakai register
              \item macro tidak dapat memakai parameter
              \item macro selalu lebih kecil
          \end{enumerate}
\end{enumerate}

\subsection{Soal Uraian}

\begin{enumerate}
    \item Jelaskan pengertian stack dan sifat LIFO. Telusuri nilai
          \reg{SP} serta isi stack untuk \lstasm|PUSH AX| dan
          \lstasm|POP BX| bila \reg{SP} mula-mula \hex{FFFE}.

    \item Jelaskan fungsi \asm{PROC}, \asm{ENDP}, \asm{CALL}, dan
          \asm{RET}. Mengapa \asm{CALL} berbeda dari \asm{JMP}?

    \item Jelaskan cara mengakses parameter melalui \eng{stack frame}
          dengan \reg{BP}. Sebutkan isi \code{[BP+0]}, \code{[BP+2]},
          \code{[BP+4]}, dan \code{[BP+6]}, lalu jelaskan fungsi
          \asm{RET} \code{4}.

    \item Jelaskan cara mendefinisikan macro berparameter dan apa
          yang terjadi saat macro dipakai. Sebutkan empat perbedaan
          prosedur dan macro.
\end{enumerate}

\subsection{Kisi-kisi Penilaian}

\begin{table}[H]
    \centering
    \caption{Kisi-kisi evaluasi Bab~\ref{chap:prosedur}.}
    \label{tab:kisi-bab14}
    \small
    \begin{tabular}{@{}p{1.5cm}p{5.0cm}p{3.0cm}c@{}}
        \toprule
        \textbf{Sub-CPMK} & \textbf{Indikator yang diukur} &
        \textbf{Bentuk soal} & \textbf{Bobot} \\
        \midrule
        3.5 & Menjelaskan stack dan PUSH/POP &
        Pilihan ganda 1; Uraian 1 & 25\% \\
        \addlinespace
        3.5 & Menjelaskan prosedur, CALL, dan RET &
        Pilihan ganda 2; Uraian 2 & 25\% \\
        \addlinespace
        3.5 & Mengakses parameter melalui stack frame &
        Pilihan ganda 3; Uraian 3 & 25\% \\
        \addlinespace
        3.5 & Mendefinisikan macro dan membedakannya &
        Pilihan ganda 4, 5; Uraian 4 & 25\% \\
        \bottomrule
    \end{tabular}
\end{table}

\begin{checklistbox}
\vspace{1.1em}
\textbf{Checklist Pencapaian Kompetensi
Bab~\ref{chap:prosedur}}

Seorang mahasiswa dinyatakan kompeten pada Sub-CPMK 3.5 apabila:
\begin{itemize}
    \item[\checked] dapat menjelaskan sifat LIFO pada stack;
    \item[\checked] dapat menggunakan PUSH dan POP;
    \item[\checked] dapat menjelaskan perubahan SP;
    \item[\checked] dapat menulis prosedur dengan PROC/ENDP;
    \item[\checked] dapat menggunakan CALL dan RET;
    \item[\checked] dapat mengakses parameter melalui BP;
    \item[\checked] dapat membuat macro berparameter;
    \item[\checked] dapat membedakan prosedur dan macro.
\end{itemize}
\end{checklistbox}


\end{document}
